\subsection{Alkalmazásarchitektúra}

Az alkalmazás felépítése több, egymástól jól elkülöníthető komponensből áll,
Ez a felépítés lehetővé teszi, hogy a felhasználói felület, állapotkezelés,
adatkezelés, adattároslás és a platformfüggő fukciók ne keveredjenek és az egyes
részek külön tesztelhetők és később könnyen bővíthetők legyenek.

\subsubsection{Legfelső réteg: Felhasználói felület}

Az alkalmazás lefelső rétegét a felhasználói felület (UI) alkotja. Ide tartoznak azok a képernők,
amelyekkel a felhasználó közvetlenül interakcióba lép. Az \code{App.kt} felel az alkalmazás
fő szerkezetéért és a képenyők közti navigációért. . A \code{ConfigScreen.kt} a mérés
előkészítéséhez kapcsolódó nézeteket tartalmazza, például a rehabilitációs feladatok kiválasztását.
Az \code{ExecutionScreen.kt} a munkamenet közbeni felületet valósítja meg, ahol a mérés vezérlése
és az aktuális gyakorlat követése történik. A \code{ResultsScreen.kt} az elkészült mérés
eredményeit jeleníti meg, míg a \code{SavedSessionScreen.kt} a korábban rögzített munkamenetek
kezeléséhez kapcsolódik. Ezek a képernyők nem önállóan végzik az adatfeldolgozást,
hanem az alkalmazás állapotkezelő és adatkezelő rétegeivel kommunikálnak.

\subsubsection{UI és háttérlogika közötti kapcsolat}

A felhasználói felület és a háttérlogika közötti kapcsolatot a \code{SessionViewModel.kt}
biztosítja. A ViewModel szerepe, hogy összefogja a munkamenet aktuális állapotát, és közvetítő
rétegként működjön a képernyők, az adattárolás, a szenzorkezelés és az elemző modulok között.
Ennek köszönhetően a képernyőknek nem kell közvetlenül ismerniük például a fájlkezelés vagy
a Bluetooth-kommunikáció részleteit.
A ViewModel kezeli, hogy a mérés milyen állapotban van, milyen adatok érkeztek be, mikor kell 
menteni az eredményeket, illetve mikor kell frissíteni a felhasználói felületen megjelenő
információkat.

\subsubsection{Adaszerkezet}

Az alkalmazás adatszerkezete külön objektumokkal van reprezentálva. A \code{SessionInfo.kt} a
munkamenethez kapcsololódó adatokat tartalmazza: a mérés azosításához és visszakereséséhez
szükséges adatokat. Az \code{ExportedTasks.kt} az adott feladatait írja le, benne van a feladat
neve és végrehajtási ideje. Az \code{ExportedSensorData} az adott nyomás és gyorsulás értékeket
gyűjti feladatoként és időlépésenként. Illetve a \code{SessionExport} rakja ezeket egybe,
tartalmazza az adott munkafolyamat információit és a gyűjtött szenzoradatokat. A
\code{SensoriaData} pedig a zokniból származó feldolgozott mérési adatokat írja le.
A \code{FullSessionResults.kt} és az \code{AnalysisModel.kt} az elemzés után létrejövő
eredmények és összesített mutatók tárolására szolgálnak. Mivel a tesztek során nem mindig
állt rendelkezésre a zokni, illetve nem volt szükség a Bluetooth kapcsolatra az adott funkciók
írásához és teszteléshez, ezért egy "hamis adat generátor" is implementálva lett az alkalmazás
írása közben ezt a logikát tartalmazza a \code{SensoriaData.kt}, a valósághoz valamennyire
közel álló szenzoradatokat hoz létre, amelyekkel a többi folyamat működése vizsgálható.
Az adatmodellek célja, hogy az alkalmazás különböző részei ugyanazokat az információkat egységes,
jól definiált formában használják.

\subsubsection{Adattárolás}

Ezen réteg feladata a mentés, betöltés és exportálás. A \code{FileStorage.kt} felelős a fájlok
létrehozásáért, olvasásáért és írásáért, amelyet a \code{SessionViewModel.kt} irányít.

\subsubsection{Szenzoros réteg}

A réteg a Sensoria okoszoknival történő kommunikációt és az onnan érkező adatok értelmezését
valósítja meg. A \code{SensoriaBluetoothManager.kt} kezeli a Bluetooth-kapcsolatot, vagyis az
eszközhöz való csatlakozást és az adatcsomagok fogadását. A beérkező nyers adatok önmagukban még
nem használhatók közvetlenül az alkalmazás többi része számára, mert hexadecimális kódként érkeznek
ezért szükség van dekódolásra. Ezt a feladatot a \code{SensoriaDecoder.kt} végzi, amely a nyers
bájtokból értelmezhető mérési értékeket állít elő. A \code{BitReader.kt} az alacsonyabb szintű
bit- és bájtolvasási műveleteket segíti, amelyek a szenzoradatok pontos feldolgozásához szükségesek. A feldolgozott értékek végül a SensoriaData.kt adatmodellben jelennek meg, így a rendszer többi része már strukturált formában tud velük dolgozni.

\subsubsection{Kamera és videó}

Ezek a funkciók szintén külön vannak kezelve. A \code{CameraPreview.kt} a mérés közbeni
kameraelőnézet megjelenítéséért felel, ezzel a felhasználó láthatja a saját mozgását a
feladat végzése közben. A \code{VideoPlayer.kt} a gyakorlatok videós megjelenítésért felelős.
Mivel ezek az osztályok operációsrendszer-specifikusak, ezért meg vannak írva Android-specifikusan is.

\subsubsection{Adatelemzés}

Az elemzési réteg a rögzített adatokból számított eredmények előállításáért felel.
A \code{SessionAnalyzer.kt} a munkamenet egészének kiértékelését végzi, vagyis a rendelkezésre
álló szenzoradatokból és feladatadatokból állít elő összesített eredményeket.
A \code{StabilityCalculator.kt} speciálisabb számításokat tartalmaz, amelyek a
stabilitással kapcsolatos mutatók meghatározására szolgálnak.
Ezek az eredmények nem diagnosztikai következtetések, hanem a munkamenet későbbi
áttekintését és összehasonlítását segítő információk.




